\chapter{Introduction}%
\label{ch:intro}%

\nv{%
In this section, I'll try to portray my ``elevator pitch''.

\textbf{Pitch:}
\begin{enumerate}
    \item \label{item:pitch1} What is the problem this thesis is trying to solve?
        \begin{itemize}
            \item The hardware and system architecture is getting more powerful.
            \item This increases the complexity of the designed system and increases the error probability.
            \item Especially since more and more junior developers can use this powerful and complicated hardware for their designs.
        \end{itemize}
    \item Why is the problem an important one to analyse?
        \begin{itemize}
            \item This can introduce huge risks in safety critical systems, risking lives and damage to humans
            \item Can cause major damage on infrastructure and commercial systems
            \item Security issues are introduced, leaking private data
        \end{itemize}
    \item Give a positive and bold statement of what this thesis is doing.
        \begin{itemize}
            \item We present novel theoretical and practically applicable analytical tools to help analyse fault-prone systems
            \item We focus on making these tools easy to use, lowering the threshold for developers to use them in everyday life
            \item \textbf{If nothing else is remembered from the thesis, this should be it.}
        \end{itemize}
    \item The consequences of the thesis (preferably ``solve'' the problem in~\ref{item:pitch1})
        \begin{itemize}
            \item By applying the tools presented in this thesis to embedded real-time control systems, we can identify and mitigate a lot of trivial bugs and errors that occur.
        \end{itemize}
\end{enumerate}
}%


%\nv{%
%\textbf{The What:}
%\begin{itemize}
%    \item Connectivity, Distributed systems, NCS, IoT, Wireless
%    \item Communication hiccups
%    \item packet drops
%    \item deadline misses
%\end{itemize}
%}%
%
%\nv{%
%\textbf{The why:}
%\begin{itemize}
%    \item Huge costs, damages
%    \item Safety
%    \item Security
%    \item Paper by~\cite{Akesson:2020}
%\end{itemize}
%}%
%
%\nv{%
%\textbf{The Statement:}
%\begin{itemize}
%    \item Theoretical and practical tools for analysing and synthesising
%        fault-tolerant controllers, with significant performance gains.
%    \item Lowering the entry-level knowledge for applying the tools to real
%        processes.
%\end{itemize}
%}%
%
%\nv{%
%\textbf{The Consequences:}
%\begin{itemize}
%    \item Decoupling and clarifying real-time and control theoretical models.
%    \item Making them more accessible to a broader class of practitioners.
%\end{itemize}
%}%

%%% WHAT
% Introduction to digitalisation
Entering the digital age has forever changed how we interact with the world and how it interacts with us.
Unlike only 20 years ago, from the moment we wake up in the morning till the moment we close our eyes at night, we interact with advanced computer systems.
Our cellphones, work computers, and even our cars going contain many computational devices, performing everything from menial tasks, such as checking the weather and accessing mail clients, to safety critical tasks, such as the car's ABS breaks and most of the engine's functionality.
Just to put the digital growth rate in perspective, the semiconductor\footnote{Semiconductors are components constituting the foundation of generally \emph{all} electronic devices.} market share has more than quadrupled over the last 20 years~\cite{statista:2022}.

% The hardware architecture is getting more powerful and cheaper to use
Not only is the number of computational devices increasing, but their independent capabilities, functionalities, and complexities are growing steadily, while the cost to buy and manufacture them has become cheaper.
Obviously, the increased efficiency and reduced cost opened up new businesses and domains, in particular within the IT-domain, whilst also consolidating and automating preexisting industry.
The integration of software in generally any domain help monitor systems, log and transmit important data, orchestrate the execution of different parts, and remotely micromanage system updates.
When everything behaves as intended, software integration is a powerful tool that improve both efficiency and increase revenue.

\note{Should I explain what type of issues I'm working with already here?}

% The complexity and reduced prices increase error surfaces
Interconnecting multiple systems is, however, not a trivial task.
When the systems are getting increasingly more complex, the surface for possible errors is also growing.
In addition to the components' individual faults, we can after connecting two components encounter problems with either the coupling or even introduce problems in the individual components.
Imagine a motorbike; it can experience all the component issues that can arise on a normal bike (punctured tire, loose chain, etc.), but it can also experience issues for connecting the bike together with a motor (electric clutch, battery, sensor networks, etc.).
Similarly to the motorbike, systems relying on data transmissions can easily be delayed or stall indefinitely if the data is lost, an orchestrator can get overloaded, and systems with remote updates have the potential to break every time a new patch is installed.
These problems are neither easy to detect nor troubleshoot; particularly since their origin can be obfuscated by complex software and hardware interconnections.

%%% WHY
% Costs
The effects of such errors can be extremely expensive and in certain cases cause companies to lose billions of dollars, even though such cases may be rare.
Typically the outcome of system faults is that the normal operation of the device (or machine) is degraded.
The degradation can over time accumulate and either wear down the device or affect the end product.
Obviously, if it is possible to extend a device's lifetime by fixing software-induced errors, a lot of money could be gained.
Furthermore, if the end product is inferior to the promised product, the consumers would go elsewhere --- no matter whether the product is a service, such as cloud storage, or a physical product, such as a cellphone.

% Safety and security
Arguably more important than the economic consequences are the risks to personal safety, security, and privacy.
One of the modern era's most famous examples are the Boeing's 737 MAX\footnote{\url{https://en.wikipedia.org/wiki/Boeing_737_MAX_groundings}} crashes, killing 346 people in two subsequent crashes; the crashes were caused by erroneous sensor readings being misinterpreted by the flight control system, resulting in the planes nosediving into the ground.
Another relevant (although less lethal) example is the famous malicious Stuxnet worm\footnote{\url{https://en.wikipedia.org/wiki/Stuxnet}}; it infiltrated the system controlling the gas-centrifuges in multiple Iranian nuclear plants, significantly damaging them whilst also collecting critical information. 
Generally it is impossible to guarantee that today's complex computer systems are absolutely safe, secure and performs according to specifications under all conditions.
Additionally, testing for all possible future problems is expensive and time consuming at best and infeasible in practice.
It is therefore crucial to develop easy-to-use, powerful tools to simplify the analysis of both the systems' performance and safety properties.

%%% STATEMENT & CONSEQUENCES
% Concluding paragraph summarising what the thesis is tackling and what the consequences of this might be
It is the purpose of this thesis to provide tools and methods for analysing systems experiencing faults.
In particular, the focus is to analyse software integrated systems where the faults are occurring in the interconnection between software and hardware.
By treating accessibility, clarity, and generalisability as first-class citizens we aim to lower the threshold for using the powerful tools provided.

\nv{Add something more to the end here?}



\subsection*{Real-Time Control Systems \fix{Need a header that better fits}}%
%
Fundamentally all systems today contain a certain level of automation, whether it is automatic heat control in buildings or memory allocation in the cloud for storing photos.
Abstracting a system to a common mathematical language and making it behave according to predefined specifications is called \emph{automatic control}.
A specific class of automatic control systems (\emph{control systems} for short) are real-time control systems, which are defined by being time critical.
The word ``real-time'' is often used synonymously with ``being very fast'', however, a more accurate synonym might be ``timely''.
In fact, a real-time system is not required to be fast, but it is required to execute according to some temporal constraints.
Figure~\ref{fig:rt-sys} depict a general setup of a real-time control system

\begin{figure}[t]
    \centering
    \def \delta {0.15}

\begin{tikzpicture}
\tikzstyle{task} = [draw,thick,fill=white,align=center]

%%% TASKS %%%

\node[task,opacity=0.3] (t1) at (-1.5+0*\delta,1.6-0*\delta) {\textcolor{white}{Task $\#3$} \\\textcolor{white}{\faFileCode[regular]}};
\node[task,opacity=0.6] (t2) at (-1.5+1*\delta,1.6-1*\delta) {\textcolor{white}{Task $\#2$} \\\textcolor{white}{\faFileCode[regular]}};
\node[task,opacity=1.0] (t3) at (-1.5+2*\delta,1.6-2*\delta) {Task $\#1$ \\\faFileCode[regular]};

\node[task,opacity=0.3] (ct1) at (1+0*\delta,1.6-0*\delta) {\textcolor{white}{Control-Task $\#3$} \\\textcolor{white}{\faFileCode[regular]}};
\node[task,opacity=0.6] (ct2) at (1+1*\delta,1.6-1*\delta) {\textcolor{white}{Control-Task $\#2$} \\\textcolor{white}{\faFileCode[regular]}};
\node[task,opacity=1.0] (ct3) at (1+2*\delta,1.6-2*\delta) {Control-Task $\#1$ \\\faFileCode[regular]};

%%% CYBER %%%

\node[thick, align=center] (rtos) at (-0.1,0.25) {Real-Time Operating System};
\node[thick, draw, align=center, rotate=90, text width=2.75cm] (hw) at (3.15,0.87) {HW Interfaces};
\node[thick, fit=(rtos)(t1)(ct1)(ct3),draw,yshift=1.5mm,xshift=0.75mm] (sw) {};
\node[thick, draw, above left] (clock) at (sw.south east) {\faClock[regular]};
\node[thick, fit=(sw)(hw), inner sep=5pt, draw] (board) {};
\node[thick, above left, xshift=1.8cm] (borad-label) at (board.south west) {Board};
\node[thick, draw, above right] (clockboard) at (board.south west)  {\faClock[regular]};

%%% PHYSICAL %%%

\node[thick, draw ,align=center] (phys) at (6,0.87) {\includegraphics[scale=4]{\figdir/airplane.jpg}};
\node[thick, draw, above left] (time) at (phys.south east) {\faClock[regular]};

%%% ARROWS %%%

\draw[-latex] ([yshift=0.65cm]hw.south) to node[yshift=0.85cm,rotate=90] {Actuation} ([yshift=0.65cm]phys.west);
\draw[-latex] ([yshift=-0.65cm]phys.west) to node[yshift=-0.85cm, rotate=90] {Sensing} ([yshift=-0.65cm]hw.south);

\end{tikzpicture}
%
    \caption{\fix{Add caption, and maybe change font? Change ``Board'' to something else? Change ``HW Interfaces to something else?''}}%
    \label{fig:rt-sys}%
\end{figure}

\nv{%
Figure (somewhere) of a general embedded setup with a brief introduction to the different
parts of the figure (nothing technical).

Short introduction to what has historically been done to mitigate the problems
}%


\subsection*{Outline}%
%
Thesis outline.
